\documentclass[12pt,a4paper]{article}

\usepackage{fontspec}
\usepackage{polyglossia}
\setmainlanguage{vietnamese}
\setmainfont{TeX Gyre Termes} % <= Sẵn trong TeX Live, tương đương Times

\usepackage{amsmath,amssymb}
\usepackage{geometry}
\usepackage{fancyhdr}
\usepackage{titlesec}
\usepackage{enumitem}
\usepackage{hyperref}
\usepackage{graphicx}
\usepackage{xcolor}
\usepackage{tabularx}
\usepackage[most]{tcolorbox}
\tcbset{
  colframe=black!40!black,
  colback=gray!2,
  boxrule=0.4pt,
  arc=1pt,
  left=4pt,
  right=4pt,
  top=2pt,
  bottom=2pt,
}

% Watermark settings
\usepackage{draftwatermark}
\SetWatermarkText{TMath Coding}     % nội dung watermark
\SetWatermarkScale{0.75}             % độ lớn (tùy chỉnh)
\SetWatermarkColor[gray]{0.92}       % màu xám nhạt (0.9 = mờ)
% End watermark settings

\geometry{top=2.5cm, bottom=2.5cm, left=2.5cm, right=2.5cm}
\setlength{\parskip}{6pt}
\setlength{\parindent}{0pt}
\setlength{\headheight}{15pt}
\pagestyle{fancy}
\fancyhf{}
\rhead{Xác suất thống kê}
\lhead{\today}
\rfoot{\thepage}

\titleformat{\section}{\large\bfseries}{\thesection.}{0.5em}{}

\begin{document}
\section*{Chương 1: Các khái niệm cơ bản}
\textbf{Biến cố ngẫu nhiên:} Một biến số mà giá trị của nó phụ thuộc vào kết quả của một hiện tượng ngẫu nhiên. Ví dụ: số mặt xuất hiện khi tung một con xúc xắc.

\textbf{Biến cố xung khắc:} Hai biến cố không thể xảy ra cùng một lúc. Ví dụ: khi tung một con xúc xắc, biến cố "xuất hiện số chẵn" và "xuất hiện số lẻ" là xung khắc.

\textbf{Biến cố độc lập:} Hai biến cố mà sự xảy ra của một không ảnh hưởng đến xác suất xảy ra của biến cố kia. Ví dụ: khi tung hai con xúc xắc, kết quả của con xúc xắc nhất không ảnh hưởng đến kết quả của con xúc xắc hai.

\textbf{Hệ đầy đủ các biến cố:} Dãy các biến cố $A_1, A_2, \ldots, A_n$ sao cho:
\begin{itemize}
    \item Các biến cố $A_i$ là xung khắc với nhau: $A_i \cap A_j = \emptyset$ với mọi $i \neq j$.
    \item Tổng của chúng là một biến cố chắc chắn: $\bigcup_{i=1}^{n} A_i = S$.
\end{itemize}

\textbf{Các tính chất của xác suất:}
\begin{itemize}
    \item $0 \leq P(A) \leq 1$ với mọi biến cố $A$.
    \item $P(\emptyset) = 0$, trong đó $\emptyset$ là biến cố không thể xảy ra.
    \item $P(S) = 1$, trong đó $S$ là biến cố chắc chắn.
    \item Nếu $A$ và $B$ là hai biến cố xung khắc, thì $P(A \cup B) = P(A) + P(B)$.
    \item Nếu $A$ và $B$ là hai biến cố bất kỳ, thì $P(A \cup B) = P(A) + P(B) - P(A \cap B)$.
    \item Với mọi biến cố $A$, ta có $P(\overline{A}) = 1 - P(A)$, trong đó $\overline{A}$ là biến cố đối của $A$.
\end{itemize}

\textbf{Xác suất điều kiện:} Xác suất của biến cố $A$ xảy ra khi biết rằng biến cố $B$ đã xảy ra, được ký hiệu là $P(A|B)$ và được tính bằng công thức:
\[P(A|B) = \frac{P(A.B)}{P(B)}, \quad P(B) > 0\]

\textbf{Công thức xác suất toàn phần:} Nếu $A_1, A_2, \ldots, A_n$ là một hệ đầy đủ các biến cố và $B$ là một biến cố bất kỳ, thì:
\[P(B) = \sum_{i=1}^{n} P(B|A_i)P(A_i)\]

\textbf{Định lý Bayes:} Nếu $A_1, A_2, \ldots, A_n$ là một hệ đầy đủ các biến cố và $B$ là một biến cố bất kỳ với $P(B) > 0$, thì:
\[P(A_i|B) = \frac{P(B|A_i)P(A_i)}{\sum_{j=1}^{n} P(B|A_j)P(A_j)}, \quad i = 1, 2, \ldots, n\]

\textbf{Dãy phép thử Bernoulli:} Một dãy các phép thử độc lập, mỗi phép thử chỉ có hai kết quả: thành công (với xác suất $p$) và thất bại (với xác suất $1-p$). Gọi $H_k$ là biến cố có đúng $k$ thành công trong $n$ phép thử, thì xác suất của biến cố này được tính bằng công thức:
\[P_n(k;p)=P(H_k) = \binom{n}{k} p^k (1-p)^{n-k}, \quad k = 0, 1, 2, \ldots, n\]

\newpage
\section*{Chương 2: Biến ngẫu nhiên và các đại lượng mô tả}
\textbf{Biến ngẫu nhiên:} Một biến số mà giá trị của nó phụ thuộc vào kết quả của một hiện tượng ngẫu nhiên.

\textbf{Biến ngẫu nhiên rời rạc:} Biến ngẫu nhiên mà nó chỉ nhận một số hữu hạn hoặc vô hạn đếm được các giá trị. Nghĩa là có thể liệt kê thành dãy $x_1, x_2, x_3, \ldots$.

\textbf{Biến ngẫu nhiên liên tục:} Biến ngẫu nhiên mà giá trị của nó có thể lấp đầy một hoặc một vài khoảng trên trục số thực và $P(X = a)$ bằng không với mọi số thực $a$.

\textbf{Hàm phân bố xác suất (CDF):} Là hàm $F(x)$ cho biết xác suất để biến ngẫu nhiên $X$ nhận giá trị nhỏ hơn hoặc bằng $x$:
\[F(x) = P(X \leq x)\]

\textbf{Tính chất của hàm phân bố xác suất:}
\begin{itemize}
    \item $0 \le F(x) \le 1$ với mọi $x \in \mathbb{R}$.
    \item $F(x)$ là hàm không giảm, liên tục bên phải. Nếu $X$ là biến ngẫu nhiên liên tục thì $F(x)$ là hàm liên tục.
    \item $F(-\infty) = 0$ và $F(+\infty) = 1$.
    \item $P(a \le x < b) = F(b) - F(a)$ với mọi $a < b$.
\end{itemize}

\textbf{Bảng phân bố xác suất:} của biến ngẫu nhiên rời rạc $X$ có dạng:
\begin{center}
\begin{tabular}{|c|c|c|c|c|c|}
\hline
$X$ & $x_1$ & $x_2$ & $x_3$ & $\ldots$ & $x_n$ \\
\hline
$P$ & $p_1$ & $p_2$ & $p_3$ & $\ldots$ & $p_n$ \\
\hline
\end{tabular}
\end{center}

\end{document}
